\chapter{Absence Seizures}
\index{Absence Seizures}

\section*{Definition and Overview}
An absence seizure is a brief episode in which a person abruptly stops what they are doing, stares blankly, and becomes unresponsive for a few seconds before recovering completely. Absence seizures are a form of generalised seizure, meaning they involve the whole brain simultaneously rather than a single region, and they were formerly known as petit mal seizures. They most commonly begin in childhood, are not accompanied by the convulsions of other seizure types, and are generally not harmful in themselves. When a child has repeated absence seizures, they may be diagnosed with childhood absence epilepsy, one of the most manageable epilepsy syndromes.

\section*{Epidemiology}
Absence seizures account for roughly 10--20\% of childhood epilepsies and affect about one or two children per 1,000. They typically begin between the ages of 4 and 10 years, are slightly more common in girls, and usually resolve by the teenage years. A family history of epilepsy increases the risk. Less commonly, absence seizures begin in adolescence or adulthood; in these people the seizures tend to persist for longer and are classified as juvenile absence epilepsy.

\section*{Aetiology and Pathophysiology}
Absence seizures arise from brief, abnormal, rhythmic bursts of electrical activity across both hemispheres of the brain at once. This synchronised discharge is thought to originate in an abnormal interaction between the cortex and the deep structures that normally regulate the brain's rhythms, particularly the thalamus. The cause is largely genetic: several genes have been associated with childhood absence epilepsy, and the underlying brain structure is typically normal. Certain triggers can provoke an episode, including hyperventilation, sleep deprivation, stress, and strong emotion.

\section*{Clinical Features}
The episodes are distinctive and often recognised first by teachers rather than parents.

\begin{itemize}
    \item Sudden cessation of activity in the middle of a task, conversation, or play
    \item A blank, staring expression and failure to respond to the child's name
    \item Episodes lasting about 5--15 seconds, sometimes with fluttering of the eyelids or slight lip or hand movements
    \item No fall, no convulsion, and no confusion afterwards; the child resumes precisely where they left off
    \item Episodes may occur many times a day, sometimes dozens of times
    \item The behaviour is often mistaken for daydreaming, inattention, or deliberately ignoring people
    \item Rapid deep breathing (hyperventilation) frequently provokes an episode
\end{itemize}

\section*{Differential Diagnosis}
A range of conditions can mimic absence seizures and must be distinguished from them.

\begin{itemize}
    \item Daydreaming, inattention, or "spacing out" in children
    \item Focal seizures with impaired awareness, which usually last longer and have a different pattern
    \item Myoclonic seizures, characterised by quick jerks rather than staring
    \item Juvenile absence epilepsy, which begins in the teenage years and tends to persist
    \item Other childhood epilepsy syndromes, such as benign rolandic epilepsy
    \item Narcolepsy and other disorders of sleep and wakefulness
    \item Drug effects or other causes of brief altered awareness
\end{itemize}

\section*{When to Seek Medical Care}
A doctor should be consulted if absence seizures are suspected, since they are often noticed first at school and a description of the episodes plus an EEG is needed for diagnosis.

\textbf{Contact your local emergency services for:}
\begin{itemize}
    \item A seizure lasting more than 5 minutes, or repeated seizures without recovery in between
    \item A first seizure with injury, difficulty breathing, or prolonged loss of consciousness
    \item A seizure occurring in water
    \item A seizure in a person with diabetes or in a pregnant woman
\end{itemize}

\section*{Diagnosis and Investigations}
The diagnosis is clinical, supported by characteristic electroencephalographic findings.

\begin{itemize}
    \item \textbf{History:} a careful eyewitness account of the episodes, including when they occur, what they look like, their duration, and what happens afterwards
    \item \textbf{Provocation by hyperventilation:} asking the child to blow or breathe rapidly can reproduce an episode in the clinic
    \item \textbf{EEG (electroencephalogram):} demonstrates the characteristic generalised spike-and-wave discharges and helps confirm the diagnosis
    \item \textbf{Video EEG:} records episodes on camera simultaneously with brain waves, providing definitive confirmation when needed
\end{itemize}

\section*{Management}
Treatment is usually medication, supported by education and safety planning.

\begin{itemize}
    \item \textbf{Medication:} ethosuximide is usually the first-line choice; alternatives include sodium valproate and lamotrigine. Because some epilepsy medicines can worsen absence seizures, treatment is tailored by a specialist
    \item \textbf{No emergency treatment:} brief absence seizures do not require rescue medication, and most children simply resume their activity
    \item \textbf{School and family education:} teachers and carers are told what the episodes look like and that the child is not being rude or inattentive
    \item \textbf{Safety advice:} avoid situations where a brief lapse could be dangerous, such as swimming alone or unsupervised cycling
    \item \textbf{Regular review:} children are followed to monitor seizure control and medication side effects
    \item \textbf{Stopping treatment:} because most children outgrow absence seizures, medication may be slowly withdrawn after a period of freedom from seizures, under specialist guidance
\end{itemize}

\section*{Complications}
\begin{itemize}
    \item Accidental injury during a brief lapse of awareness, including falls and traffic incidents
    \item Problems at school when episodes are frequent or unrecognised, affecting concentration and learning
    \item Social and emotional difficulties from being labelled as inattentive or disruptive
    \item A minority of children later develop generalised tonic-clonic (grand mal) seizures
    \item Prolonged absence status, in which the child remains confused and unresponsive for longer than usual, is uncommon but possible
\end{itemize}

\section*{Prognosis and Outcomes}
The outlook for childhood absence epilepsy is excellent. The great majority of children outgrow the seizures by the teenage years, and most have normal intelligence and development. Seizures that begin in adolescence are more likely to continue into adulthood and usually require longer-term treatment. With appropriate medication, most children achieve complete control of their seizures and can expect a full and active quality of life.

\section*{Prevention and Self-Care}
\begin{itemize}
    \item Give medication regularly and do not stop it without specialist advice
    \item Maintain good sleep routines, since tiredness can trigger seizures
    \item Avoid excessive stress and overheating
    \item Ensure school staff know what the episodes look like and what to do
    \item Supervise swimming and bathing, and use helmets where appropriate when cycling
    \item Attend regular reviews with the treating doctor or neurologist
\end{itemize}

\section*{Sources and Further Reading}
The following organisations provide reliable information on absence seizures and childhood epilepsy.

\begin{itemize}
    \item NHS. \textit{Absence seizures: overview and treatment.} https://www.nhs.uk/conditions/absence-seizures/
    \item Mayo Clinic. \textit{Petit mal seizure: symptoms and causes.} https://www.mayoclinic.org/diseases-conditions/petit-mal-seizure/
    \item MSD Manual. \textit{Seizure disorders: professional reference.} https://www.msdmanuals.com/
    \item MedlinePlus. \textit{Absence seizure health information.} https://medlineplus.gov/ency/article/000696.htm
    \item Epilepsy Foundation. \textit{Absence seizures information for families.} https://www.epilepsy.com/
    \item NICE. \textit{Epilepsies: diagnosis and management.} https://www.nice.org.uk/guidance/ng217
\end{itemize}
